\par
Sea $\underline A$ un grupoide cofibrado homogéneo sobre
$\underline C_\Lambda$, y sea $f:R'\longrightarrow R$ una pequeña
extensión en $\underline C_\Lambda$. Existe una aplicación canónica

%%%%%%%%%%  scan idx 67  |  folio 56  %%%%%%%%%%

\[
 \widetilde f:R'\times_RR'\longrightarrow k[\operatorname{Ker}f],
 \qquad
 (r'_1,r'_2)\longmapsto r'_1(0)+(r'_1-r'_2),
\]
y
\[
 1\times\widetilde f:R'\times_RR'
 \longrightarrow R'\times_k k[\operatorname{Ker}f]
\]
es un isomorfismo. Dados dos $f$-morfismos
$a_i\longrightarrow a$ $(i=1,2)$ en $\underline A$, el producto
fibrado $a_1\times_a a_2$ es un objeto sobre $R'\times_RR'$. Por
tanto, podemos formar
\[
 \tau(a_1,a_2)=
 (\widetilde f)_*(a_1\times_a a_2)
 \in\underline A(k[\operatorname{Ker}f]).
\]
Como $1\times\widetilde f$ es un isomorfismo, existe un
isomorfismo
\[
 a_1\times_a a_2\simeq a_1\times\tau(a_1,a_2)
\]
compatible con la proyección sobre el primer factor. El enunciado
1.5(2) adopta ahora la forma más precisa siguiente.

\underline{Lema 2.8.} Sea $\underline A$ un grupoide cofibrado
homogéneo sobre $\underline C_\Lambda$, y sean
\[
 a_i\xrightarrow{\ \alpha_i\ }a\qquad(i=1,2)
\]
$f$-morfismos en $\underline A$.

(1) Existe $\rho:a_1\longrightarrow a_2$ en $\underline A$
con $\alpha_2\circ\rho=\alpha_1$ si y solo si existe una flecha
$a_1\longrightarrow\tau(a_1,a_2)$ en $\underline A$.

(2) Existe tal $\rho$ con $P(\rho)=\operatorname{id}_{R'}$
si y solo si
\[
 \tau(a_1,a_2)=0
 \quad\text{en}\quad
 [\underline A(k[\operatorname{Ker}f])].
\]

\underline{Demostración.} Demostramos simultáneamente ambas
afirmaciones. La existencia de $\rho:a_1\longrightarrow a_2$ que
satisface $\alpha_2\circ\rho=\alpha_1$ (y
$P(\rho)=\operatorname{id}_{R'}$) equivale sucesivamente a la
existencia de
\[
 h:a_1\longrightarrow a_1\times_a a_2,
 \qquad \operatorname{pr}_1\circ h=\operatorname{id}_{a_1},
\]
con $P(h)$ igual a la aplicación diagonal; a la existencia de
\[
 g:a_1\longrightarrow a_1\times\tau(a_1,a_2),
 \qquad \operatorname{pr}_1\circ g=\operatorname{id}_{a_1},
\]
con
\[
 P(g):R'\longrightarrow R'\times_k k[\operatorname{Ker}f],
 \qquad r'\longmapsto(r',r'(0));
\]
y a la existencia de una flecha
$\gamma:a_1\longrightarrow\tau(a_1,a_2)$ en $\underline A$, con
\[
 P(\gamma):R'\longrightarrow k[\operatorname{Ker}f],
 \qquad r'\longmapsto r'(0).
\]
Tal flecha se factoriza a través de
$R'\longrightarrow k\longrightarrow k[\operatorname{Ker}f]$ si y
solo si $\tau(a_1,a_2)=0$. \hfill$\square$

\underline{Proposición 2.9.} (a) Sean
\[
 \underline A\xrightarrow{\ \varphi\ }\underline B
 \xrightarrow{\ \psi\ }\underline C
\]
funtores de grupoides cofibrados sobre $\underline C_\Lambda$.
Supongamos que

(i) $\underline B$ es semihomogéneo y $\underline C$ es
homogéneo;

(ii) la aplicación
\[
 \psi(k[\varepsilon]):[\underline B(k[\varepsilon])]
 \longrightarrow[\underline C(k[\varepsilon])]
\]
es inyectiva.

Entonces $\psi\circ\varphi$ es liso si y solo si $\varphi$ y
$\psi$ son ambos

%%%%%%%%%%  scan idx 68  |  folio 57  %%%%%%%%%%

lisos.

(b) Sea $\varphi:\underline A\longrightarrow\underline B$ un
funtor liso de grupoides cofibrados sobre $\underline C_\Lambda$.
Si $\underline A$ es semihomogéneo y $\underline B$ es
homogéneo, entonces
\[
 \widehat\varphi:\widehat{\underline A}
 \longrightarrow\widehat{\underline B}
\]
es liso. Explícitamente, para toda aplicación sobreyectiva
$b'\xrightarrow{\beta}\widehat\varphi(a)$ en
$\widehat{\underline B}$, existen una aplicación
$\alpha:a'\longrightarrow a$ en $\widehat{\underline A}$ y un
$\widehat{\underline C}_\Lambda$-isomorfismo
\[
 \rho:\widehat\varphi(a')\xrightarrow{\ \sim\ }b'
\]
tales que $\beta\circ\rho=\varphi(\alpha)$.

\underline{Demostración.} (a) Supongamos que
$\psi\circ\varphi$ es liso. Si $\varphi$ es liso, entonces
$\psi$ es liso por la Proposición 2.3(b), de modo que basta
demostrar que $\varphi$ es liso. Sea
$b'\xrightarrow{\beta}\varphi(a)$ una pequeña prolongación en
$\underline B$. Como $\psi\circ\varphi$ es liso, existen una
flecha $a'\xrightarrow{\alpha}a$ en $\underline A$ y un
$\underline C_\Lambda$-isomorfismo
\[
 \rho:\psi(b')\longrightarrow\psi\varphi(a')
\]
para los cuales el diagrama
\[
\begin{tikzcd}[column sep=huge,row sep=huge]
 \psi(b') \arrow[rr,"{\rho}"] \arrow[dr,"{\psi(\beta)}"']
 && \psi\varphi(a') \arrow[dl,"{\psi\varphi(\alpha)}"]\\
 & \psi\varphi(a) &
\end{tikzcd}
\]
es conmutativo. Como $\underline B$ es semihomogéneo, podemos
elegir un diagrama conmutativo
\[
\begin{tikzcd}[column sep=small,row sep=small]
 & b'' \arrow[dl] \arrow[dr] &
 &&& R'\times_RR' \arrow[dl] \arrow[dr] &\\
 b' \arrow[dr,"{\beta}"'] &&
 \varphi(a') \arrow[dl,"{\varphi(\alpha)}"]
 & \text{situado sobre} &
 R' \arrow[dr] && R' \arrow[dl]\\
 & \varphi(a) &&&& R &
\end{tikzcd}
\]
donde $R=P(a)$ y $R'=P(b')=P(a')$. La homogeneidad de
$\underline C$ da
\[
 \psi(b'')=
 \psi(b')\times_{\psi\varphi(a)}\psi\varphi(a')
\]
por la Observación 2.6(a). El Lema 2.8(2) da entonces
\[
 0=\tau(\psi\varphi(a'),\psi(b'))
   =(\widetilde f)_*\psi(b'').
\]
Por la hipótesis (ii), se sigue que $(\widetilde f)_*(b'')=0$.
El enunciado 1.5(2) proporciona ahora una aplicación
$\gamma:b'\longrightarrow\varphi(a')$ que satisface
$\beta=\varphi(\alpha)\circ\gamma$. Por tanto, $\varphi$ es liso,
y entonces también lo es $\psi$. El recíproco se deduce de la
Proposición 2.3(a).

(b) Sea dada una aplicación sobreyectiva
$\beta:b'\longrightarrow\widehat\varphi(a)$ en
$\widehat{\underline B}$. Está representada por una sucesión de
diagramas conmutativos

%%%%%%%%%%  scan idx 69  |  folio 58  %%%%%%%%%%

\[
\begin{tikzcd}[column sep=large]
 {} \arrow[r,dashed] & b'_3 \arrow[r,"{\beta'_2}"] \arrow[d,"{\beta_3}"]
 & b'_2 \arrow[r,"{\beta'_1}"] \arrow[d,"{\beta_2}"]
 & b'_1 \arrow[d,"{\beta_1}"]\\
 {} \arrow[r,dashed] & \varphi(a_3) \arrow[r,"{\varphi(\alpha_2)}"']
 & \varphi(a_2) \arrow[r,"{\varphi(\alpha_1)}"']
 & \varphi(a_1)
\end{tikzcd}
\]
en $\underline B$, donde cada $\alpha_v$, $\beta_v$ y
$\beta'_v$ es sobreyectiva y
$\beta=\varprojlim_v\beta_v$. Basta demostrar por inducción que
todo diagrama de flechas de trazo continuo que aparece a
continuación, con
$\alpha_n$, $\beta_n$, $\beta_{n+1}$ y $\beta'_n$
sobreyectivas, puede completarse mediante las flechas punteadas:
\[
\begin{tikzcd}[column sep=large,row sep=large]
 \varphi(a'_{n+1})
   \arrow[rrrr,dashed,"{\varphi(\,)}"]
   \arrow[dr,dashed,leftrightarrow]
   \arrow[ddr,dashed,"{\varphi(\,)}"']
 &&&& \varphi(a'_n)
   \arrow[dl,leftrightarrow]
   \arrow[ddl,"{\varphi(\alpha'_n)}"]\\
 & b'_{n+1} \arrow[rr,"{\beta'_n}"] \arrow[d,"{\beta_{n+1}}"]
 && b'_n \arrow[d,"{\beta_n}"'] &\\
 & \varphi(a_{n+1}) \arrow[rr,"{\varphi(\alpha_n)}"']
 && \varphi(a_n) &
\end{tikzcd}
\]
Por hipótesis, $\underline A$ es semihomogéneo y
$\sisi{\alpha_n}{\alpha'_n}:a'_{n}\longrightarrow a_n$ es sobreyectiva. Por tanto,
existe un diagrama conmutativo
\[
\begin{tikzcd}[column sep=small,row sep=small]
 & a' \arrow[dl,"{\gamma}"'] \arrow[dr,"{\gamma'}"] &
 &&& R_{n+1}\times_{R_n}R'_n \arrow[dl] \arrow[dr] &\\
 a_{n+1} \arrow[dr,"{\alpha_n}"'] &&
 a'_n \arrow[dl,"{\alpha'_n}"]
 & \text{en $\underline A$ sobre} &
 R_{n+1} \arrow[dr] && R'_n \arrow[dl]
 & \text{en $\underline C_\Lambda$}\\
 & a_n &&&& R_n &&
\end{tikzcd}
\]
La homogeneidad de $\underline B$ implica que
\[
\begin{tikzcd}[column sep=huge,row sep=huge]
 & \varphi(a') \arrow[dl,"{\varphi(\gamma)}"']
                  \arrow[dr,"{\varphi(\gamma')}"] &\\
 \varphi(a_{n+1}) \arrow[dr,"{\varphi(\alpha_n)}"']
 && \varphi(a'_n) \arrow[dl,"{\varphi(\alpha'_n)}"]\\
 & \varphi(a_n) &
\end{tikzcd}
\]

%%%%%%%%%%  scan idx 70  |  folio 59  %%%%%%%%%%

es un diagrama de producto fibrado en $\underline B$. Existe, por
tanto, una única aplicación
$\delta:b'_{n+1}\longrightarrow\varphi(a')$ con la conmutatividad
requerida. Como $\beta_{n+1}$ y $\beta'_n$ son sobreyectivas,
también lo es $\delta$. La lisura de $\varphi$ proporciona entonces
una aplicación $a'_{n+1}\xrightarrow{\alpha'}a'$ en $\underline A$
y un $\underline C_\Lambda$-isomorfismo
\[
 \rho:\varphi(a'_{n+1})\xrightarrow{\ \sim\ }b'_{n+1}
\]
tal que $\delta\circ\rho=\varphi(\alpha')$. Obtenemos el diagrama
conmutativo
\[
\begin{tikzcd}[row sep=huge,column sep=huge]
 \varphi(a'_{n+1})
   \arrow[rrrr,"{\varphi(\gamma'\circ\alpha')}"]
   \arrow[dr,"{\rho}"]
   \arrow[ddr,"{\varphi(\gamma\circ\alpha')}"']
 &&&& \varphi(a'_n)
   \arrow[dl,"{\sim}"]
   \arrow[ddl,"{\varphi(\alpha'_n)}"]\\
 & b'_{n+1} \arrow[rr] \arrow[d]
 && b'_n \arrow[d] &\\
 & \varphi(a_{n+1}) \arrow[rr]
 && \varphi(a_n) &
\end{tikzcd}
\]
lo que completa la inducción y demuestra (b). \hfill$\square$
